% Document layout: body font, one-time font measurement, uniform leading and
% page setup.  This ties together the geometry of the pitch marks (pitch-marks.tex)
% and the held-syllable underline (hold.tex), so it must be input *after* both.

% Helvetica at 18pt.  The Adobe metric phvr8t has a 10pt design size,
% so "scaled 1800" gives 18pt.  \risescale (set in pitch-marks.tex) sizes the
% marks: 1.8 keeps them proportional to the text, and it is currently 0.9 for
% half size.
\font\rm=phvr8t scaled 1800
\rm
% Measure the body font once.  \riseheight (its cap/ascender height) fixes how
% far above the baseline every pitch mark sits; the depth feeds the leading.
\setbox0=\hbox{Ahg}%
\riseheight=\ht0
\dimen1=\dp0                              % descender depth
\holddepth=\dimen1 \advance\holddepth by1.5pt    % \hold underline: this far below baseline
\holdthick=1.44pt                                % ... and this thick (independent of \risescale)
\holdgap=4pt                                     % ... and this much narrower than the syllable
% Uniform leading: give every line the same height -- with or without marks --
% for an even vertical rhythm.  Set it once from the mark geometry and force
% TeX to always use \baselineskip (\lineskiplimit=-\maxdimen), so every pair of
% baselines is exactly \baselineskip apart regardless of a line's contents.
% (This replaces the per-mark strut that used to be added in \risemark.)
\dimen2=8bp \dimen2=\risescale\dimen2     % strokes rise this far above the raise point
\advance\dimen2 by\riseheight             % \dimen2 = top of the marks above the baseline
\topskip=\dimen2 \advance\topskip by2pt   % first line: room for its marks + small margin
\advance\dimen2 by\holddepth              % + room below for a \hold underline...
\advance\dimen2 by\holdthick              %   (down to its lower edge)
\advance\dimen2 by3pt                     % + a small gap between lines
\baselineskip=\dimen2
\lineskiplimit=-\maxdimen                 % always honour \baselineskip => identical lines
\parindent=0pt
% Hanging indent: the first line of each paragraph stays flush left (parindent
% is 0); every following (wrapped) line is indented.  \hangafter=1 applies
% \hangindent to the lines after the first.  Both reset each paragraph, so set
% them for every paragraph (including those starting with a mark) via \everypar.
\everypar={\hangindent=1.5em \hangafter=1 }
\pdfpagewidth=210mm \pdfpageheight=297mm   % A4 paper
% Symmetric horizontal margins.  Plain TeX's page origin is 1in from the left
% and \hoffset is 0, so the left margin is 1in; setting \hsize = pagewidth - 2in
% makes the right margin 1in too.  Left, right and top margins are then all
% 1 inch (25.4mm).  (The default \hsize of 6.5in is a US-Letter width, which on
% narrower A4 paper left an asymmetric ~19.5mm right margin.)
\hsize=\pdfpagewidth \advance\hsize by-2in
